\section{Satisfaction}
\label{sec:fol-semantics-02}

\begin{definition}[Satisfaction]
For a structure \(\Meta\), assignment \(s\), and formula \(\varphi\), the
relation \(\Meta,s\entails\varphi\) is defined by structural recursion:
\[
\begin{aligned}
 \Meta,s\entails r=t
 &\quad\text{iff}\quad
 \sem{r}^{\Meta}_s=\sem{t}^{\Meta}_s,\\
 \Meta,s\entails R(t_1,\ldots,t_n)
 &\quad\text{iff}\quad
 (\sem{t_1}^{\Meta}_s,\ldots,\sem{t_n}^{\Meta}_s)\in R^{\Meta},\\
 \Meta,s\notentails\bottom,
 &\\
 \Meta,s\entails\neg\varphi
 &\quad\text{iff}\quad\Meta,s\notentails\varphi,\\
 \Meta,s\entails\varphi\land\psi
 &\quad\text{iff}\quad
 \Meta,s\entails\varphi\text{ and }\Meta,s\entails\psi,\\
 \Meta,s\entails\varphi\lor\psi
 &\quad\text{iff}\quad
 \Meta,s\entails\varphi\text{ or }\Meta,s\entails\psi,\\
 \Meta,s\entails\varphi\to\psi
 &\quad\text{iff}\quad
 \Meta,s\notentails\varphi\text{ or }\Meta,s\entails\psi,\\
 \Meta,s\entails\forall x\,\varphi
 &\quad\text{iff}\quad
 \Meta,s[x\mapsto a]\entails\varphi
 \text{ for every }a\in\dom\Meta,\\
 \Meta,s\entails\exists x\,\varphi
 &\quad\text{iff}\quad
 \Meta,s[x\mapsto a]\entails\varphi
 \text{ for some }a\in\dom\Meta.
\end{aligned}
\]
\end{definition}

\begin{lemma}[Formula coincidence]
\label{lem:fol-coincidence}
If assignments \(s\) and \(s'\) agree on \(\FV(\varphi)\), then
\[
 \Meta,s\entails\varphi
 \quad\text{iff}\quad
 \Meta,s'\entails\varphi.
\]
\end{lemma}
\begin{proof}
Proceed by structural induction.  Atomic cases use term coincidence.
Connective cases follow from the induction hypotheses.  For
\(\forall x\,\psi\), the assignments
\(s[x\mapsto a]\) and \(s'[x\mapsto a]\) agree on \(\FV(\psi)\) for each
\(a\), because \(s\) and \(s'\) agree on
\(\FV(\psi)\setminus\set{x}\).  Apply the induction hypothesis for every
\(a\).  The existential case is identical with ``some'' in place of
``every.''
\end{proof}

\begin{theorem}[Substitution lemma]
\label{thm:fol-substitution}
For every formula \(\varphi\), term \(t\), variable \(x\), structure
\(\Meta\), and assignment \(s\),
\[
 \Meta,s\entails\varphi[t/x]
 \quad\text{iff}\quad
 \Meta,s[x\mapsto\sem{t}^{\Meta}_s]\entails\varphi.
\]
\end{theorem}
\begin{proof}
Proceed by induction on formula complexity.  Atomic cases follow from
\cref{lem:term-substitution}.  Connective cases follow immediately from the
induction hypotheses.

Let \(\varphi=Qy\,\psi\), where
\(Q\in\set{\forall,\exists}\).  If \(y=x\), substitution stops at the
quantifier.  Updating \(x\) before evaluating \(Qx\,\psi\) is immaterial,
because the quantified update of \(x\) overrides it.  If \(y\ne x\) and
\(y\notin\operatorname{Var}(t)\), then
\(\sem{t}^{\Meta}_{s[y\mapsto a]}=\sem{t}^{\Meta}_s\) for every \(a\).
Apply the induction hypothesis to \(\psi\) under each updated assignment and
then apply the truth clause for \(Q\).

If \(y\ne x\) and \(y\in\operatorname{Var}(t)\), the definition first replaces \(y\) by a
fresh \(z\).  The variable \(z\) occurs neither in \(t\) nor free in
\(\psi\).  The induction hypothesis applied to the renamed body---
legitimate because renaming preserves complexity, by
\cref{lem:substitution-complexity}---together with formula coincidence,
shows directly that
\(Qy\,\psi\) and \(Qz\,(\psi[z/y])\) have the same satisfaction condition.
The preceding quantifier argument then applies to the renamed body, whose
complexity equals \(c(\psi)\) by \cref{lem:substitution-complexity}.  Thus
the equivalence holds in all cases.
\end{proof}

\begin{corollary}[Alpha-invariance]
\label{cor:alpha-invariance}
If \(\varphi\equiv_\alpha\psi\), then
\(\Meta,s\entails\varphi\) iff \(\Meta,s\entails\psi\) for every structure
\(\Meta\) and assignment \(s\).
\end{corollary}
\begin{proof}
Define a relation \(R\) on formulas by \(\varphi R\psi\) iff, for every
structure \(\Meta\) and assignment \(s\),
\(\Meta,s\entails\varphi\) iff \(\Meta,s\entails\psi\).  The relation
\(R\) is reflexive, symmetric, and transitive.  It is a congruence for
\(\bottom,\neg,\land,\lor,\to\) by the corresponding satisfaction clauses.
It is also a congruence for each quantifier: if \(\varphi R\psi\), then the
equivalence holds under every assignment \(s[x\mapsto a]\), and the
universal or existential satisfaction clause transfers it to
\(Qx\,\varphi\) and \(Qx\,\psi\).

It remains to check a generating renaming.  Suppose
\(y\notin\operatorname{Var}(\varphi)\).  By
\cref{thm:fol-substitution}, satisfaction of \(\varphi[y/x]\) under
\(s[y\mapsto a]\) is equivalent to satisfaction of \(\varphi\) under
\(s[y\mapsto a][x\mapsto a]\).  Since \(y\notin\FV(\varphi)\), formula
coincidence identifies the latter satisfaction condition with that of
\(\varphi\) under \(s[x\mapsto a]\).  Applying the satisfaction clause for
\(Q\in\set{\forall,\exists}\) gives
\(Qx\,\varphi\;R\;Qy\,(\varphi[y/x])\).  Thus \(R\) is an equivalence
relation and a congruence containing every generating renaming.  By the
leastness clause in the definition of \(\equiv_\alpha\), alpha-equivalence is
contained in \(R\).
\end{proof}

\begin{corollary}[Assignment independence for sentences]
If \(\sigma\in\Sent(\Lang)\), then
\(\Meta,s\entails\sigma\) iff \(\Meta,s'\entails\sigma\) for all assignments
\(s,s'\) in \(\Meta\).
\end{corollary}
\begin{proof}
The assignments agree on \(\FV(\sigma)=\varnothing\); apply
\cref{lem:fol-coincidence}.
\end{proof}

\begin{definition}[Truth, validity, and semantic consequence]
For a sentence \(\sigma\), write \(\Meta\entails\sigma\) iff
\(\Meta,s\entails\sigma\) for one, equivalently every, assignment \(s\).
For \(\Gamma\subseteq\Fm(\Lang)\), write
\(\Gamma\entails\varphi\) iff, for every \(\Lang\)-structure \(\Meta\) and
assignment \(s\),
\[
 \bigl(\Meta,s\entails\gamma\text{ for every }\gamma\in\Gamma\bigr)
 \quad\Longrightarrow\quad
 \Meta,s\entails\varphi.
\]
The formula \(\varphi\) is \emph{valid} iff
\(\varnothing\entails\varphi\).
Two \(\Lang\)-formulas \(\varphi\) and \(\psi\) are \emph{logically
equivalent}, written \(\varphi\equiv\psi\), iff, for every
\(\Lang\)-structure \(\Meta\) and every assignment
\(s\),
\[
  \Meta,s\entails\varphi
  \quad\Longleftrightarrow\quad
  \Meta,s\entails\psi.
\]
This is logical equivalence over all \(\Lang\)-structures, not
alpha-equivalence and not literal identity of formula strings.
\end{definition}

